\documentclass[11pt]{article}
\usepackage[margin=0.85in]{geometry}
\usepackage{booktabs,hyperref}
\hypersetup{colorlinks=true,urlcolor=blue}
\title{Impala performance experiment: Step 1}
\author{Cache fixed material-model preparation}
\date{September 29, 2026}
\begin{document}\maketitle
\section*{Change and scope}
The material evaluator now reuses fixed experimental conditions, integration
histories, strain offsets, and interpolation coordinates. A four-entry least
recently used cache checks array contents, batch size, and history resolution.
In-place edits invalidate the cache. Copies protect cached arrays from mutable
physics state and preserve the original memory layout. Actual simulated strain
coordinates are checked before reuse.

Material parameters and constants are evaluated fresh. Physical equations,
parameter validation, SciPy linear interpolation, 100-step histories, floating
point precision, and MCMC algorithms are unchanged. This is a CPU preparation
optimization; no GPU kernel is introduced. The normal material-model grid setter
is assumed; externally stateful custom grid setters are outside this cache's
supported inputs. A model instance remains mutable and is not thread safe.
\section*{Accuracy evidence}
All 66 tests passed in 100.93 seconds, including 14 new regressions for cache
invalidation, array mutation, batch changes, old instances, exceptional inputs,
and numerical agreement. The independent original fixture contains 256 parameter
sets: 128 feasible prior inputs and 128 sampled inputs. Predictions and weighted
log likelihoods matched exactly, with maximum absolute error zero.

Four paired sampler experiments (pooled and hierarchical, seeds 42 and 123)
each used 3,000 stored states, eight temperatures, six copper experiments and
327 observations. Original and cached versions produced exactly equal parameter,
noise, and acceptance arrays at every temperature. Hierarchical means,
covariances and swap counters also matched exactly. These tests compare the
frozen original evaluator against the optimized evaluator with shared physics
and samplers; future physics changes require a separate original checkout too.
\section*{Measured speed}
Model timings are medians of 60 paired repetitions, alternating execution order.
Compilation is warmed before measurement. Speedup is original time divided by
cached time. These measurements are specific to this machine and workload.
\begin{center}\begin{tabular}{rrrr}\toprule
Batch & Original (ms) & Cached (ms) & Speedup\\\midrule
1 & 1.545 & 1.501 & 1.029\\
4 & 1.798 & 1.727 & 1.041\\
50 & 5.247 & 4.816 & 1.090\\
200 & 15.889 & 14.359 & 1.107\\
\bottomrule\end{tabular}\end{center}
Complete sampler timings include the same inference work on both sides; execution
order is reversed for seed 123.
\begin{center}\begin{tabular}{lrrrr}\toprule
Mode & Seed & Original (s) & Cached (s) & Speedup\\\midrule
pooled & 42 & 7.721 & 7.527 & 1.026\\
pooled & 123 & 8.092 & 7.762 & 1.043\\
hier & 42 & 11.700 & 10.319 & 1.134\\
hier & 123 & 10.787 & 10.381 & 1.039\\
\bottomrule\end{tabular}\end{center}
The improvement is modest; the main physics history loop still dominates.
Two seeds per mode are insufficient to establish a precise general speedup.
Short-chain equality is a strong regression check, not a convergence claim.
No long-run ESS-per-second study was repeated for this patch. The previously
observed hierarchical mixing problem remains unresolved.
\section*{Tracking and reproduction}
Original baseline: \texttt{6ccd5a1}, tagged \texttt{perf-baseline-20260929}.\\
Initial implementation: \texttt{7ac9310}. Final measured code: \texttt{5c5cb01}.\\
Branch: \texttt{perf/cache-material-preparation}.\\
Final results checkpoint: \texttt{perf-step1-20260929} (local tag).

Versioned documentation is under \texttt{docs/performance/}: the change log,
accuracy gates, compact results, test logs, artifact hashes, and this LaTeX
source. Full traces and raw timings are preserved outside the checkout in
\texttt{benchmark/changes/step1-final/}. Earlier measurements remain in a separate
\texttt{step1/} directory. The original baseline PDF is preserved.

Run ordinary \texttt{pytest} for regression coverage; the existing pull-request
workflow collects these tests. Run the paired benchmark explicitly:
\begin{verbatim}
.venv/bin/python benchmarks/cache_material_preparation.py \
  --output ../benchmark/changes/step1-rerun
\end{verbatim}
Choose a new output directory each time. Never regenerate original fixtures
using optimized code. Every later optimization should have its own commit,
accuracy results, measured timings, and change-log entry. If arithmetic changes
prevent exact trajectories, assess numerical tolerances and multi-seed posterior
and predictive agreement, R-hat, ESS, and ESS per second before accepting it.
\end{document}
